% ══════════════════════════════════════════════════════════════════════════
% §3  The Library Layer: BLAS & Eigen
% ══════════════════════════════════════════════════════════════════════════
\section{The Library Layer: BLAS \& Eigen}
\label{sec:blas}

\subsection{Why Libraries Beat Hand-Written Loops}

The \textbf{Basic Linear Algebra Subprograms (BLAS)} specification defines
three levels of operations:

\begin{center}
\renewcommand{\arraystretch}{1.3}
\begin{tabular}{@{}clcc@{}}
\toprule
\textbf{Level} & \textbf{Operation} & \textbf{FLOPs} & \textbf{Data} \\
\midrule
1 & $\alpha x + y$ \quad (AXPY)       & $O(n)$   & $O(n)$ \\
2 & $Ax + y$ \quad (GEMV)             & $O(n^2)$ & $O(n^2)$ \\
3 & $\alpha AB + \beta C$ \quad (GEMM) & $O(n^3)$ & $O(n^2)$ \\
\bottomrule
\end{tabular}
\end{center}

\begin{tipbox}
Level~3 (GEMM) is the ``Gold Standard'' because its arithmetic intensity
is $O(n^3)/O(n^2) = O(n)$: the CPU reuses each cache line many times before
evicting it.  This is why matrix--matrix multiply runs close to peak FLOP/s,
while matrix--vector multiply is memory-bound.
\end{tipbox}

Optimised BLAS implementations---\textbf{Intel MKL}, \textbf{OpenBLAS},
\textbf{BLIS}---apply cache-blocking, SIMD, and multi-threading
automatically.  A na\"ive triple loop achieves $<5\%$ of peak;
a tuned GEMM reaches $>90\%$.

\subsection{The SpMV Speed Trap}

Engineers often assume that because they are using a ``fast'' library like
Eigen or MKL, their sparse solver will hit peak FLOP/s.
\textbf{It will not.}

When you solve a sparse system iteratively (e.g.\ GMRES on the
$10^6 \times 10^6$ heat rod), the inner loop is a \emph{Sparse
Matrix--Vector multiply} (SpMV)---a Level~2 operation.  Its arithmetic
intensity is $O(\text{nnz})/O(\text{nnz}) = O(1)$: each non-zero is
loaded from RAM, multiplied once, and never reused.

\begin{warningbox}
SpMV is \textbf{memory-bandwidth bound}, not compute-bound.
No amount of SIMD or MKL magic can fix a slow RAM bus if you are
just streaming a million doubles once per iteration.  On a typical
desktop with 40~GB/s bandwidth, SpMV on a tridiagonal $10^6$ system
peaks at $\sim\!3$~GFLOP/s---less than $5\%$ of the CPU's
theoretical throughput.
\end{warningbox}

This is why \emph{preconditioning} (\S\ref{sec:gmres}) matters so much:
reducing the iteration count by $100\times$ is worth far more than any
micro-optimisation of the SpMV kernel itself.

\subsection{Eigen as a BLAS Wrapper}

When you write
\begin{lstlisting}
Eigen::MatrixXd C = A * B;
\end{lstlisting}
Eigen's \emph{expression templates} detect this is a GEMM and dispatch to
whichever BLAS backend is linked at compile time.  The cascade is:

\begin{enumerate}[nosep]
    \item \textbf{MKL linked?} $\rightarrow$ calls \texttt{cblas\_dgemm}
          (fastest on Intel hardware).
    \item \textbf{OpenBLAS linked?} $\rightarrow$ calls its GEMM
          (good cross-platform default).
    \item \textbf{Neither?} $\rightarrow$ Eigen's built-in blocked GEMM
          (decent, but slower than the above).
\end{enumerate}

\begin{warningbox}
If you install Eigen but never link a BLAS library, you are leaving
a $2$--$5\times$ speedup on the table for every dense multiply.
\end{warningbox}

\subsection{Linking MKL in Practice}

With CMake and Intel oneAPI installed:
\begin{lstlisting}[language=bash,morekeywords={find_package,target_link_libraries}]
# CMakeLists.txt
find_package(MKL REQUIRED)
target_link_libraries(my_app PRIVATE MKL::MKL)
add_definitions(-DEIGEN_USE_MKL_ALL)
\end{lstlisting}

The following snippet benchmarks na\"ive loops against Eigen's GEMM:

\lstinputlisting[caption={BLAS-level GEMM: na\"ive loop vs.\ Eigen.},
                 label={lst:gemm}]
                {code/gemm_benchmark.cpp}
